\documentclass[12pt,table,xcolor={dvipsnames}]{beamer}
\usetheme{Pittsburgh}
\usecolortheme{seagull}
%\usepackage[utf8]{inputenc}
\usepackage{fontspec}
\usepackage[brazilian]{babel}
\usepackage{amsmath}
\usepackage{listings}
\usepackage{multirow}
\usepackage{amsfonts}
\usepackage{amssymb}
\usepackage{graphicx}
\author{Programação C++}
\title{Estruturas de Dados}
%\setbeamercovered{transparent}
\setbeamertemplate{navigation symbols}{}
%\logo{\includegraphics[scale=0.015]{Brasao_UFSC.png}\includegraphics[scale=0.2]{brasao_PPGCC.jpg}}
\institute{Departamento de Informática e de Estatística \\ Prof. Jean Everson Martina \\ Prof. Aldo von Wangenheim}
\date{2016.2}
\subject{}
\usebackgroundtemplate{\includegraphics[width=\paperwidth,
height=\paperheight]{../reusable_images/fundo_UFSC.png}}
\begin{document}

{
\usebackgroundtemplate{\includegraphics[width=\paperwidth,
height=\paperheight]{../reusable_images/fundo_capa.png}}
\begin{frame}
\titlepage
\includegraphics[scale=0.3]{../reusable_images/brasao_INE.png}
\end{frame}
}

\begin{frame}{Classes em C++}
\begin{itemize}
\item Declaração das classes é semelhante ao Java (Java é que copiou do C++);
\item As melhores práticas são declarar no arquivo Header e implementar no Source;
\item Isso tem por objetivo separar instrumentação do compilador de código
objeto.
\end{itemize}
\end{frame}

\begin{frame}[fragile]{\textit{Pessoa.h}}
\lstset{language=C++,
          keywordstyle=\color{blue}\ttfamily,
          stringstyle=\color{red}\ttfamily,
          commentstyle=\color{OliveGreen}\ttfamily,
          breaklines=true,
          basicstyle=\ttfamily\footnotesize
          }
\begin{lstlisting}
class Pessoa {
 private:
  // declaracao dos atributos privados
 protected:
  // declaracao dos metodos protegidos
 public:
  // declaracao dos metodos publicos
}
\end{lstlisting}
\end{frame}

\begin{frame}[fragile]{\textit{Pessoa.cpp}}
\lstset{language=C++,
          keywordstyle=\color{blue}\ttfamily,
          stringstyle=\color{red}\ttfamily,
          commentstyle=\color{OliveGreen}\ttfamily,
          breaklines=true,
          basicstyle=\ttfamily\footnotesize
          }
\begin{lstlisting}
#include "Pessoa.h"

void Pessoa::metodo1() {
// implementacao do metodo 1
}

int Pessoa::metodo2() {
// implementacao do metodo 2
}
\end{lstlisting}
\end{frame}

\begin{frame}{Classes em C++}
\begin{itemize}
\item Estas são as melhores práticas;
\item Nada impede de você implementar os métodos no Header;
\item Templates vão exigir a implementação dentro do header pois o pré-processador do compilador que tem que criar as variações.
\end{itemize}
\end{frame}

\begin{frame}[fragile]{Construtores e Destrutores}
\lstset{language=C++,
          keywordstyle=\color{blue}\ttfamily,
          stringstyle=\color{red}\ttfamily,
          commentstyle=\color{OliveGreen}\ttfamily,
          breaklines=true,
          basicstyle=\ttfamily\footnotesize
          }
\begin{lstlisting}
// Pessoa.h
class Pessoa {
 private:
  // declaracao dos atributos privados
 protected:
  // declaracao dos metodos protegidos
 public:
  Pessoa(); // construtor padrao
  Pessoa(char n[], int i);//  com parametros
  virtual ~Pessoa(); // destrutor
}
\end{lstlisting}
\end{frame}

\begin{frame}[fragile]{Construtores e Destrutores}
\lstset{language=C++,
          keywordstyle=\color{blue}\ttfamily,
          stringstyle=\color{red}\ttfamily,
          commentstyle=\color{OliveGreen}\ttfamily,
          breaklines=true,
          basicstyle=\ttfamily\footnotesize
          }
\begin{lstlisting}
// Pessoa.cpp
#include "Pessoa.h"

Pessoa::Pessoa() {
// implementacao
}

int Pessoa::Pessoa(char* n, int i) {
// implementacao
}

int Pessoa::~Pessoa() {
// implementacao
}
\end{lstlisting}
\end{frame}

\begin{frame}[fragile]{Construtores e Destrutores}
\lstset{language=C++,
          keywordstyle=\color{blue}\ttfamily,
          stringstyle=\color{red}\ttfamily,
          commentstyle=\color{OliveGreen}\ttfamily,
          breaklines=true,
          basicstyle=\ttfamily\footnotesize
          }
Usando os construtores e destrutores:
\begin{lstlisting}
Pessoa p; // construtor padrao e invocado
Pessoa p("Fulano", 27);//outro construtor
p.~Pessoa(); // destrutor
Pessoa* p = new Pessoa(); //dinamico
Pessoa* p = new Pessoa("Fulano", 27); //dinamico
delete p; //dinamico
\end{lstlisting}
\end{frame}

\begin{frame}[fragile]{Herança em C++}
\lstset{language=C++,
          keywordstyle=\color{blue}\ttfamily,
          stringstyle=\color{red}\ttfamily,
          commentstyle=\color{OliveGreen}\ttfamily,
          breaklines=true,
          basicstyle=\ttfamily\footnotesize
          }
\begin{lstlisting}
Class Empregado : public Pessoa {
// herda atributos e metodos da superclasse igual em java
}
Class Empregado : protected Pessoa {
// herda atributos e metodos da superclasse com visao so pra classe e para seus filhos
}
Class Empregado : private Pessoa {
// herda atributos e metodos da superclasse com visao somente interna
}
\end{lstlisting}
\end{frame}

\begin{frame}[fragile]{Métodos Virtuais}
\lstset{language=C++,
          keywordstyle=\color{blue}\ttfamily,
          stringstyle=\color{red}\ttfamily,
          commentstyle=\color{OliveGreen}\ttfamily,
          breaklines=true,
          basicstyle=\ttfamily\footnotesize
          }
\begin{lstlisting}
Class Empregado : public Pessoa {
  private:
   int _salario;
   int _descontos;
  public:
   virtual int salarioLiquido() = 0;
}
\end{lstlisting}
\end{frame}

\begin{frame}[fragile]{Métodos Virtuais}
\lstset{language=C++,
          keywordstyle=\color{blue}\ttfamily,
          stringstyle=\color{red}\ttfamily,
          commentstyle=\color{OliveGreen}\ttfamily,
          breaklines=true,
          basicstyle=\ttfamily\footnotesize
          }
\begin{lstlisting}
Class Estudante : public Empregado {
  private:
  public:
   virtual int salarioLiquido() {
    return _salario;
   };
   virtual void adicionaBonus() {
    _salario = _salario + 0;
   }
}
\end{lstlisting}
\end{frame}

\begin{frame}[fragile]{Métodos Virtuais}
\lstset{language=C++,
          keywordstyle=\color{blue}\ttfamily,
          stringstyle=\color{red}\ttfamily,
          commentstyle=\color{OliveGreen}\ttfamily,
          breaklines=true,
          basicstyle=\ttfamily\footnotesize
          }
\begin{lstlisting}
Class Professor : public Estudante {
  private:
  public:
   int salarioLiquido() {
    return _salario * 0.725;
   };
   void adicionaBonus() {
    _salario = _salario + 5000;
   }
}
\end{lstlisting}
\end{frame}

\begin{frame}{Classes Abstratas e Interfaces}
\begin{itemize}
\item Classe abstrata tem pelo menos um método puramente virtual (sem implementação);
\item Interface tem todos os métodos puramente virtuais;
\item Em C++trudo se dá por herança múltipla;
\item Isso causa problemas de ambiguidade como o Problema do Diamante;
\item A Solução é fazer herança virtual:
\begin{itemize}
\item public virtual;
\item protected virtual;
\item private virtual.
\end{itemize}
\end{itemize}
\end{frame}

\begin{frame}{Tratamento de Exceções}
\begin{itemize}
\item O tratamento de exceções permite capturar erros ocorridos durante a execução de um programa;
\item Uma exceção é lançada com o comando \textit{throw()};
\item O comando \textit{throw()} deve ser usado dentro de um bloco \textit{try\{...\}};
\item O fluxo de execução é desviado para o bloco \textit{catch()}.
\end{itemize}
\end{frame}

\begin{frame}[fragile]{Tratamento de Exceções}
\lstset{language=C++,
          keywordstyle=\color{blue}\ttfamily,
          stringstyle=\color{red}\ttfamily,
          commentstyle=\color{OliveGreen}\ttfamily,
          breaklines=true,
          basicstyle=\ttfamily\footnotesize
          }
\begin{lstlisting}
// exceptions
#include <iostream>
using namespace std;

int main () {
  try
  {
    throw 20;
  }
  catch (int e)
  {
    cout << "An exception occurred. Exception Nr. " << e << '\n';
  }
  return 0;
}
\end{lstlisting}
\end{frame}

\begin{frame}[fragile]{Exceções Padronizadas no C++}
\lstset{language=C++,
          keywordstyle=\color{blue}\ttfamily,
          stringstyle=\color{red}\ttfamily,
          commentstyle=\color{OliveGreen}\ttfamily,
          breaklines=true,
          basicstyle=\ttfamily\footnotesize
          }
\begin{lstlisting}
// using standard exceptions
#include <iostream>
#include <exception>
using namespace std;

class myexception: public exception {
  virtual const char* what() const throw() {
    return "My exception happened";
  }
} myex;

int main () {
  try {
    throw myex;
  } catch (exception& e) {
    cout << e.what() << '\n';
  }
  return 0;
}
\end{lstlisting}
\end{frame}

\begin{frame}{Templates}
\begin{itemize}
\item A utilização de templates nos permite criar funções e classes genéricas;
\item As funções ou classes genéricas são baseadas em argumentos com tipo desconhecido;
\item O tipo será definido quando forem usadas por entidades específicas.
\end{itemize}
\end{frame}

\begin{frame}[fragile]{Templates}
\lstset{language=C++,
          keywordstyle=\color{blue}\ttfamily,
          stringstyle=\color{red}\ttfamily,
          commentstyle=\color{OliveGreen}\ttfamily,
          breaklines=true,
          basicstyle=\ttfamily\footnotesize
          }
\begin{lstlisting}
template<typename T>
class Lista {
 private:
  Lista *proximo;
  T *dado;
 public:
  T retornaDoInicio();
};
int main() {
Lista<int> listaInteiro;
Lista<char*> listaString;
}
\end{lstlisting}
\end{frame}

\begin{frame}{Documentação com Doxygen}
\begin{itemize}
\item  Ele gerar  documentação de referência a partir de um conjunto de códigos documentados;
\item  Tem suporte para gerar saída em RTF (Word), PostScript, PDF com links, HTML compactado e man pages do Linux.;
\item A documentação é extraída diretamente dos códigos;
\item Você pode também gerar automaticamente as relações entre os vários elementos, o que inclui gráficos de dependência, diagramas de herança e diagramas colaborativos;
\end{itemize}
\end{frame}

\begin{frame}[fragile]{Documentação com Doxygen}
\lstset{language=C++,
          keywordstyle=\color{blue}\ttfamily,
          stringstyle=\color{red}\ttfamily,
          commentstyle=\color{OliveGreen}\ttfamily,
          breaklines=true,
          basicstyle=\ttfamily\footnotesize
          }
\begin{lstlisting}
//!  Uma classe de teste.
/*!
  Uma descricao mais elaborada da classe.
*/
class Teste {
 public:
  //! Uma enumeracao.
  /*! Descricao mais detalhada da enumeracao. */
  enum TEnum {
               TVal1, /*!< Valor enumeravel TVal1. */
               TVal2, /*!< Valor enumeravel TVal2. */
               TVal3  /*!< Valor enumeravel TVal3. */
             }

\end{lstlisting}
\end{frame}

\begin{frame}[fragile]{Documentação com Doxygen}
\lstset{language=C++,
          keywordstyle=\color{blue}\ttfamily,
          stringstyle=\color{red}\ttfamily,
          commentstyle=\color{OliveGreen}\ttfamily,
          breaklines=true,
          basicstyle=\ttfamily\footnotesize
          }
\begin{lstlisting}
         //! Ponteiro da enumeracao.
         /*! Detalhes. */
         *ptrEnum,
         //! Variavel da enumeracao.
         /*! Detalhes. */
         varEnum;
    //! Um construtor.
    /*!
      Descricao mais elaborada do construtor.
    */
    Teste();

    //! Um destrutor.
    /*!
      Descricao mais elaborada do destrutor.
    */
   ~Teste();
\end{lstlisting}
\end{frame}

\begin{frame}[fragile]{Documentação com Doxygen}
\lstset{language=C++,
          keywordstyle=\color{blue}\ttfamily,
          stringstyle=\color{red}\ttfamily,
          commentstyle=\color{OliveGreen}\ttfamily,
          breaklines=true,
          basicstyle=\ttfamily\footnotesize
          }
\begin{lstlisting}
  //! Um membro normal com dois argumentos e retornando um valor inteiro.
  /*!
    \param a um argumento inteiro.
    \param s um ponteiro para char const.
    \return Os resultados do teste
    \sa Teste(), ~Teste(), meTesteTambem() and varPublica()
  */
  int meTeste(int a,const char *s);
  //! Um membro puramente virtual.
  /*!
    \sa testMe()
    \param c1 o primeiro argumento.
    \param c2 o segundo argumento.
  */
 virtual void meTesteTambem(char c1,char c2)=0;
\end{lstlisting}
\end{frame}

\begin{frame}[fragile]{Documentação com Doxygen}
\lstset{language=C++,
          keywordstyle=\color{blue}\ttfamily,
          stringstyle=\color{red}\ttfamily,
          commentstyle=\color{OliveGreen}\ttfamily,
          breaklines=true,
          basicstyle=\ttfamily\footnotesize
          }
\begin{lstlisting}
    //! Uma variavel publica.
    /*!
      Detalhes.
    */
    int varPublica;

    //! Uma variavel funcao.
    /*!
      Detalhes.
    */
    int (*handler)(int a,int b);
};
\end{lstlisting}
\end{frame}

{
\usebackgroundtemplate{\includegraphics[width=\paperwidth,
height=\paperheight]{../reusable_images/fundo_capa.png}}
\begin{frame}

{\LARGE Perguntas????}

\end{frame}
}


{
\usebackgroundtemplate{\includegraphics[width=\paperwidth,
height=\paperheight]{../reusable_images/fundo_capa.png}}
\begin{frame}
\includegraphics[scale=0.8]{../reusable_images/cc_logo_arge.png}\hspace{0.5cm}
\includegraphics[scale=0.95]{../reusable_images/by.png}

\vspace{1cm}
Este trabalho está licenciado sob uma Licença Creative Commons Atribuição 4.0 Internacional. Para ver uma cópia desta licença, visite http://creativecommons.org/licenses/by/4.0/.

\end{frame}
}
\end{document}
